\chapter[PENDAHULUAN]{\\ PENDAHULUAN}

\section{Latar Belakang Masalah}
Perkembangan teknologi informasi dan komunikasi telah mendorong terjadinya transformasi digital di berbagai sektor bisnis, termasuk dalam pengelolaan transaksi dan operasional penjualan pada usaha mikro, kecil, dan menengah (UMKM). Pemanfaatan sistem \textit{Point of Sale} (POS) berbasis aplikasi menjadi salah satu solusi yang banyak digunakan untuk membantu pelaku usaha dalam mencatat transaksi, mengelola stok barang, serta menyusun laporan penjualan secara lebih terstruktur dan efisien. Seiring dengan meningkatnya kebutuhan akan pengambilan keputusan yang cepat dan berbasis data, sistem POS tidak lagi hanya berfungsi sebagai alat pencatatan transaksi, tetapi juga dituntut untuk mampu menyediakan informasi analitis yang mendukung strategi bisnis.

Pada praktiknya, sebagian besar aplikasi \textit{Point of Sale} yang digunakan oleh UMKM masih bersifat deskriptif, yaitu hanya menampilkan data historis penjualan tanpa adanya kemampuan analisis prediktif. Pelaku usaha umumnya harus melakukan perhitungan dan estimasi penjualan secara manual berdasarkan pengalaman atau intuisi, yang berpotensi menimbulkan kesalahan dalam perencanaan stok, penentuan target penjualan, dan pengelolaan arus kas. Ketidakmampuan sistem dalam memprediksi penjualan harian secara akurat dapat menyebabkan berbagai permasalahan, seperti kelebihan stok, kekurangan stok, serta tidak optimalnya strategi penjualan yang diterapkan.

Perkembangan teknologi kecerdasan buatan, khususnya dalam bidang \textit{machine learning} dan \textit{deep learning}, membuka peluang untuk menghadirkan fitur prediksi penjualan yang lebih akurat dan adaptif terhadap pola data historis. Salah satu metode \textit{deep learning} yang banyak digunakan dalam analisis data runtun waktu (\textit{time series}) adalah \textit{Long Short-Term Memory} (LSTM). Metode LSTM memiliki keunggulan dalam mempelajari pola jangka pendek dan jangka panjang pada data penjualan harian, sehingga mampu menghasilkan prediksi yang lebih baik dibandingkan metode statistik konvensional.

Penerapan metode \textit{Long Short-Term Memory} (LSTM) dalam sistem \textit{Point of Sale} diharapkan dapat membantu pelaku usaha dalam memprediksi penjualan harian berdasarkan data transaksi sebelumnya. Dengan adanya fitur prediksi penjualan, sistem POS tidak hanya berperan sebagai alat pencatatan transaksi, tetapi juga sebagai sistem pendukung keputusan yang dapat memberikan gambaran tren penjualan di masa mendatang. Informasi prediksi tersebut dapat dimanfaatkan untuk perencanaan stok barang, pengaturan strategi promosi, serta pengelolaan keuangan usaha secara lebih efektif dan efisien.

Berdasarkan permasalahan tersebut, maka diperlukan suatu penelitian yang merancang dan membangun aplikasi \textit{Point of Sale} dengan fitur prediksi penjualan harian menggunakan metode \textit{Long Short-Term Memory} (LSTM). Penelitian ini diharapkan dapat menghasilkan sebuah aplikasi POS yang mampu memberikan nilai tambah bagi pengguna melalui pemanfaatan teknologi kecerdasan buatan, serta menjadi referensi bagi pengembangan sistem \textit{Point of Sale} berbasis \textit{machine learning} di masa mendatang.

\section{Rumusan Masalah}
Berdasarkan latar belakang yang telah diuraikan, maka rumusan masalah dalam penelitian ini adalah sebagai berikut:
\begin{enumerate}
    \item Bagaimana merancang dan membangun aplikasi \textit{Point of Sale} (POS) dengan fitur prediksi penjualan harian berbasis \textit{mobile} yang dapat digunakan oleh pelaku usaha secara efektif?
    \item Bagaimana menerapkan metode \textit{Long Short-Term Memory} (LSTM) pada aplikasi \textit{Point of Sale} untuk memprediksi penjualan harian berdasarkan data transaksi historis?
    \item Bagaimana tingkat akurasi hasil prediksi penjualan harian yang dihasilkan oleh metode \textit{Long Short-Term Memory} (LSTM) berdasarkan evaluasi menggunakan metrik kesalahan prediksi?
\end{enumerate}

\section{Tujuan Penelitian}
Berdasarkan rumusan masalah yang telah diuraikan, maka tujuan dalam penelitian ini adalah sebagai berikut:
\begin{enumerate}
    \item Merancang dan membangun aplikasi \textit{Point of Sale} (POS) berbasis \textit{mobile} yang dilengkapi dengan fitur prediksi penjualan harian sebagai alat bantu pengelolaan dan analisis penjualan.
    \item Mengimplementasikan metode \textit{Long Short-Term Memory} (LSTM) pada aplikasi \textit{Point of Sale} untuk memprediksi penjualan harian berdasarkan data transaksi historis.
    \item Memanfaatkan data transaksi penjualan harian, seperti jumlah transaksi, total penjualan, dan jumlah produk terjual, sebagai data latih utama dalam model prediksi penjualan berbasis metode LSTM.
    \item Mengevaluasi kinerja dan tingkat akurasi metode \textit{Long Short-Term Memory} (LSTM) dalam menghasilkan prediksi penjualan harian menggunakan metrik evaluasi kesalahan prediksi.
\end{enumerate}

\section{Batasan Masalah}
Untuk menghindari meluasnya materi pembahasan tugas akhir ini, maka pembahasan akan dibatasi mencakup hal-hal berikut:
\begin{enumerate}
    \item Aplikasi dikembangkan menggunakan framework Flutter dan difokuskan untuk sistem operasi Android, mengingat dominasi pasar Android pada segmen UMKM lokal.
    \item Studi kasus dilakukan pada UMKM skala mikro di wilayah Kota Langsa.
    \item Algoritma yang digunakan adalah \textit{Long Short-Term Memory} (LSTM) dengan model \textit{univariate} atau \textit{multivariate} terbatas, dan prediksi dilakukan untuk rentang waktu harian.
    \item Sistem mencakup pencatatan penjualan kasir, manajemen stok sederhana, dan laporan penjualan harian. Fitur pembayaran digital (\textit{e-wallet integration}) tidak dibahas secara mendalam pada versi ini.
    \item Proses pelatihan model (\textit{training}) dan prediksi dilakukan di sisi \textit{Server/Cloud} menggunakan bahasa Python, sedangkan aplikasi \textit{mobile} bertugas mengirim data transaksi dan menerima hasil prediksi melalui \textit{REST API}.
\end{enumerate}

\section{Manfaat Penelitian}
Manfaat yang diharapkan dari penelitian ini adalah sebagai berikut:
\begin{enumerate}
    \item \textbf{Bagi Pelaku UMKM di Kota Langsa}: Memberikan aplikasi \textit{Point of Sale} (POS) berbasis \textit{mobile} yang dapat digunakan untuk pencatatan transaksi penjualan dan pengelolaan stok secara digital, serta menyediakan fitur prediksi penjualan harian berbasis algoritma \textit{Long Short-Term Memory} (LSTM) sebagai alat bantu dalam pengambilan keputusan \textit{restock} dan pengelolaan persediaan barang.
    \item \textbf{Bagi Akademisi}: Menjadi referensi dan bahan kajian dalam pengembangan penelitian di bidang sistem informasi dan \textit{machine learning}, khususnya penerapan metode \textit{Long Short-Term Memory} (LSTM) untuk prediksi penjualan harian pada sistem \textit{Point of Sale} (POS) berbasis \textit{mobile} di lingkungan UMKM.
    \item \textbf{Bagi Peneliti}: Menambah pengetahuan dan pengalaman dalam merancang dan membangun aplikasi POS berbasis \textit{mobile} menggunakan framework Flutter serta mengimplementasikan algoritma \textit{Long Short-Term Memory} (LSTM) untuk menyelesaikan permasalahan prediksi penjualan harian berbasis data transaksi UMKM pada studi kasus nyata.
\end{enumerate}

\section{Sistematika Penulisan}
Sistematika penulisan proposal skripsi ini disusun untuk memberikan gambaran secara menyeluruh mengenai alur pembahasan yang dilakukan. Sistematika penulisan dibagi menjadi tiga bab utama yang saling berkaitan, dengan rincian sebagai berikut:

\noindent\textbf{BAB I PENDAHULUAN}\\
Bab ini menguraikan gambaran umum penelitian yang meliputi latar belakang masalah, rumusan masalah, tujuan penelitian, batasan masalah, manfaat penelitian, serta sistematika penulisan. Bab ini menjadi dasar dalam memahami alasan dan arah penelitian yang dilakukan.

\noindent\textbf{BAB II TINJAUAN PUSTAKA}\\
Bab ini berisi landasan teori dan kajian literatur yang relevan dengan penelitian. Teori yang disajikan mencakup konsep \textit{Point of Sale} (POS), peramalan (\textit{forecasting}), data deret waktu (\textit{time series}), \textit{machine learning}, konsep \textit{Recurrent Neural Network} (RNN), \textit{Long Short-Term Memory} (LSTM), serta metrik evaluasi kinerja model. Selain itu, bab ini juga menyajikan ringkasan penelitian terdahulu (\textit{State of the Art}) dalam bentuk tabel komparatif.

\noindent\textbf{BAB III METODOLOGI PENELITIAN}\\
Bab ini menjelaskan tahapan metodologi penelitian yang dilakukan, meliputi alur penelitian, pengumpulan data, perancangan sistem menggunakan UML (\textit{use case} dan \textit{activity diagram}), pemodelan menggunakan metode LSTM, rancangan pengujian sistem menggunakan metode \textit{black-box testing}, serta jadwal dan anggaran kegiatan penelitian.